% Generated draft: requires mathematical and source review.
\documentclass[11pt]{article}
\usepackage[margin=25mm]{geometry}
\usepackage{fontspec}
\setmainfont{Latin Modern Roman}
\usepackage{amsmath,amssymb,mathtools,unicode-math}
\setmathfont{Latin Modern Math}
\usepackage{enumitem,needspace,xurl,hyperref}
\setlength{\parindent}{0pt}
\setlength{\parskip}{5pt}
\setlength{\emergencystretch}{3em}
\newcommand{\N}{\mathbb{N}}
\newcommand{\Z}{\mathbb{Z}}
\newcommand{\Q}{\mathbb{Q}}
\newcommand{\R}{\mathbb{R}}
\newcommand{\C}{\mathbb{C}}
\newcommand{\F}{\mathbb{F}}
\newcommand{\A}{\mathbb{A}}
\newcommand{\PP}{\mathbb{P}}
\newcommand{\E}{\mathbb{E}}
\newcommand{\eps}{\varepsilon}
\newcommand{\dd}{\,\mathrm{d}}
\begin{document}
\section*{UnsolvedMath ID: 30003343}
\textbf{Source problem code:} catalogue entry
\subsection{Definitions and mathematical statement}
\textbf{Retained research question.}
\begin{quote}
How do historical categories or personae used to describe women in science influence our view of female mathematicians today?
\end{quote}
\textbf{Limits of mathematical formulation.} The supplied record identifies this as a historical, biographical, or sociological research prompt, rather than an open mathematical problem. It supplies no mathematical definitions, domains, hypotheses, quantifiers, or proposition to prove or disprove. In particular, it does not specify the historical categories or personae, the historical periods, whose contemporary views are under consideration, the time denoted by ``today,'' or what constitutes influence. These are unresolved aspects of the research question; no mathematical formalization can be recovered from the supplied evidence without adding substantive assumptions.

The registry marks the statement as \texttt{unrecoverable}, supplies no clean statement, and separately labels the record \texttt{open}. These are supplied registry classifications, not independently verified mathematical conclusions.
\subsection{Short English statement}
The question asks how historical descriptions of women in science shape present views of female mathematicians. The supplied record treats it as a historical or sociological inquiry. It does not provide enough structure for a precise mathematical statement; the original question is retained for human review.
\subsection{Sources}
\begin{itemize}
\item Supplied local UnsolvedMath registry record, permanent ID \(30003343\), title ``Historical Scientific Personae and Women Mathematicians,'' exact source code \texttt{catalogue entry}, in the Oberwolfach Reports Open Problems collection. No individual proposer is supplied. The recorded proposed year is 2017.
\item Supplied source citation: \emph{Mini-Workshop: Women in Mathematics: Historical and Modern Perspectives} (2018), \emph{Oberwolfach Reports}, pp.\ 83--131, DOI \texttt{10.4171/owr/2017/2}. \url{https://doi.org/10.4171/owr/2017/2}. The record describes Oberwolfach Reports as documenting mathematical workshops and their open-problem discussions. The citation and report attribution are supplied and have not been independently checked here.
\item Supplied literature source: \emph{Against All Odds: Women's Ways to Mathematical Research}. \url{https://doi.org/10.1007/978-3-030-47610-6}.
\item Supplied literature assessment, checked \texttt{2026-08-21}: status \texttt{open}. The supplied assessment says that historical studies show that gendered personae and categories shaped recognition, but do not settle how past categories determine views of contemporary female mathematicians. Its listed sources are the Oberwolfach report and \emph{Against All Odds}, linked above. This is supplied dated curation derived from the OWR record, not an independently verified literature assessment; time-sensitive status and cited claims require independent verification.
\end{itemize}
\end{document}
